% =============================================================================
%                    TEMPLATE VOOR CURSUSSAMENVATTING
%                    KU Leuven - School Macros v4.0
% =============================================================================
% Dit is een starttemplate voor nieuwe documenten.
% Kopieer dit bestand en pas het aan voor je eigen cursus.
%
% Tip: Compileer TWEEMAAL voor correcte formularium en symbolenlijst!
% =============================================================================

\documentclass[a4paper,11pt]{article}
\usepackage{../../school-macros}
\usetikzlibrary{shapes.multipart}

% v4.0 Optioneel: Compact Mode voor cheatsheets
% \enablecompactmode

% =============================================================================
%                           DOCUMENT INFO
% =============================================================================
% Pas deze gegevens aan voor je cursus:
\newcommand{\vaknaam}{Object gericht-programmeren}
\newcommand{\auteur}{Ruben Ryckaert}

\title{\vaknaam}
\subtitle{Samenvatting + synthax}
\author{\auteur}

\date{\today}


% =============================================================================
%                              DOCUMENT
% =============================================================================
\begin{document}

% dutch babel maakt " actief ("a wordt a-umlaut); hier uitzetten
\ifdefined\shorthandoff\shorthandoff{"}\fi
\maketitle


\begin{abstract}
    Deze samenvattingen zijn beschikbaar op GitHub waar je kunt bijdragen aan verbeteringen: \url{https://github.com/KUL-Industriele-ingenieurs/Samenvattingen-Ku-leuven-Industriele-ingenieurs}

    Je bent helemaal vrij en ik raad je zelf aan om info toe te voegen aan de documenten als je de uitleg niet goed vindt.
    Je leert je vak en je raakt dan ook bekend met LaTeX wat handig is voor later.
\end{abstract}

\vspace{1cm}

\subsection{Vak en examen info}

In het vak objectgericht programmeren leer je Java en leer je objectgericht redeneren. Meer info vind je op studforum \url{https://wiki.studforum.net/pub/bac2/oop}

Je mag voor het examen een \textbf{handgeschreven
    4-pagina's (2 bladen voor en achter) A4-formularium meenemen} met daarop alles wat je nodig hebt.
Dit document moet je tonen wat de belangrijkste dingen zijn die je moet weten voor het examen.
Ik raad je aan om dit document te lezen en de belangrijkste dingen over te schrijven op je formularium.
Je hebt ook op studforum grote hoeveelheden oefenexamens staan die je zeker moet maken.
\url{https://cloud.studforum.net/apps/files/?dir=/GTS-Fase2/Semester%201/Courses%20After%202021/Nederlandse%20Groep/Objectgerichte%20Softwareontwikkeling&fileid=794166}
\newline

W3schools heeft een goede Java tutorial: \url{https://www.w3schools.com/java/}
Bekijk die zeker als je niet weet wat alles doet. Dit is een goed beginpunt.

De slides zijn ook niet zoveel dus als je een herhaling nodig hebt kun je de slides ook even doorgaan.

Je moet je vooral focussen op oefeningen maken en snel kunnen redenenen en programmeren. Deze samenvatting toont alleen de synthax en alle soorten
concepten die je gaat toepassen in de oefeningen maar je moet er ook echt maken.

Gebruik ook Bluej (\url{https://www.bluej.org/}), het programma dat je \textbf{moet gebruiken op het examen}. Leer er code in te typen, \textit{ctrl + space (geeft je overview
    over alle code die je kunt gebruiken)}, en leer er mee debuggen. Meer info over Bluej in het boek en op de slides.


\subsection{De Java API-documentatie op het examen}
\label{sec:oop-javadocs}

Op het examen staat de pc in examenmodus: geen internet. Je krijgt wel een
offline kopie van de Java API-documentatie (het bestand \texttt{docs.zip}, net
zoals in de oefenzittingen). Dat is exact dezelfde documentatie als
\url{https://docs.oracle.com/en/java/javase/17/docs/api/}, maar dan lokaal.

\textbf{Dit is je belangrijkste hulpmiddel tijdens het programmeergedeelte.}
De methodenamen van klassen zoals \texttt{ArrayList} of \texttt{Scanner} hoef je
dus niet vanbuiten te kennen --- je moet ze wel \emph{kunnen terugvinden}. Oefen
dat op voorhand, want zoeken kost tijd die je op het examen niet hebt.

\begin{enumerate}
    \item Pak \texttt{docs.zip} uit en open \texttt{index.html} in de browser.
    \item Typ de klassenaam in de zoekbalk rechtsboven, bv.\ \texttt{HashMap}.
          Of ga via \textit{All Classes} en zoek met \textit{Ctrl+F}.
    \item Op de pagina van een klasse gebruik je twee tabellen:
          \begin{itemize}
              \item \textbf{Constructor Summary} --- hoe maak je zo'n object aan?
              \item \textbf{Method Summary} --- alle methodes, met per methode het
                    \textbf{returntype} (linkerkolom), de \textbf{parameters} en een
                    regel uitleg.
          \end{itemize}
    \item Lees altijd het returntype. Dat vertelt je meteen of je het resultaat
          in een variabele moet opvangen (\texttt{boolean}, \texttt{int},
          \texttt{V}\dots) of niet (\texttt{void}).
\end{enumerate}

\begin{examenbox}
    Dit zijn de klassen die je op het examen zeker moet kunnen terugvinden.
\end{examenbox}

\begin{center}
    \begin{tabular}{@{}llp{6.0cm}@{}}
        \toprule
        \textbf{Klasse}    & \textbf{Package}   & \textbf{Waarvoor}                                                       \\
        \midrule
        \texttt{ArrayList} & \texttt{java.util} & lijst die groeit en krimpt                                              \\
        \texttt{HashMap}   & \texttt{java.util} & paren sleutel--waarde                                                   \\
        \texttt{HashSet}   & \texttt{java.util} & verzameling zonder dubbels                                              \\
        \texttt{Iterator}  & \texttt{java.util} & verwijderen tijdens het doorlopen                                       \\
        \texttt{Scanner}   & \texttt{java.util} & lezen van toetsenbord of bestand                                        \\
        \texttt{Random}    & \texttt{java.util} & willekeurige getallen                                                   \\
        \texttt{String}    & \texttt{java.lang} & tekst bewerken en vergelijken                                           \\
        \texttt{Math}      & \texttt{java.lang} & \texttt{abs}, \texttt{max}, \texttt{round}, \texttt{sqrt}, \texttt{pow} \\
        \texttt{File}      & \texttt{java.io}   & een bestand aanduiden                                                   \\
        \bottomrule
    \end{tabular}
\end{center}

Alles uit \texttt{java.lang} (zoals \texttt{String} en \texttt{Math}) is altijd
beschikbaar: daarvoor heb je g\'e\'en \texttt{import} nodig. Voor de rest wel.

De namen in de kolom \textbf{Klasse} hierboven zijn wat je in de zoekbalk van
\texttt{docs.zip} typt. Verderop in deze samenvatting verwijs ik daar niet telkens
opnieuw naar: staat er een methode die je niet kent, zoek de klasse dan gewoon op
in de docs.


\subsection{Tips}

\textbf{Volg de practica.} Vier uur is lang, maar de info die je haalt uit het
\textit{fout} oplossen van een oefening --- en daar dan feedback over vragen aan de
assistent --- is essentieel om het vak onder de knie te krijgen. Je hoeft niet
altijd de volledige zitting vol te houden, maar \textbf{vanaf zitting 5 of 6}
beginnen de moeilijkere oefeningen.

\textbf{Bereid de zittingen voor} en werk voorgaande onafgewerkte zittingen af,
zodat je gerichte vragen kunt stellen aan de assistenten.

\textbf{De hoorcolleges} zijn een pak minder nuttig: je leert niet programmeren
door de les letterlijk over te nemen. De theorie achter Java is ook niet van
belang voor het examen.

\textbf{Het boek is onhandig.} Online staan veel betere cursussen die je sneller
op weg helpen (zie \textit{Nuttige links} verderop).



\subsection{Nuttige links}
\label{sec:oop-links}

\begingroup\sloppy
\begin{itemize}
    \item \textbf{Studforum wiki OOP} --- algemene vakinfo en tips: \\
          \url{https://wiki.studforum.net/pub/bac2/oop}
    \item \textbf{Oefenexamens} op de studforum-cloud: \\
          \url{https://cloud.studforum.net/apps/files/?dir=/GTS-Fase2/Semester\%201/Courses\%20After\%202021/Nederlandse\%20Groep/Objectgerichte\%20Softwareontwikkeling&fileid=794166}
    \item \textbf{W3Schools Java-tutorial} --- goed beginpunt: \\
          \url{https://www.w3schools.com/java/}
\end{itemize}
\endgroup

\section{Basis java}
\subsection{Import statements}
Om klassen van andere packages te gebruiken,
moeten deze eerst geïmporteerd worden met een import statement.\\

\begin{codeblock}[Import utils]{java}
    import java.util.*; // importeert alle klassen van de package java.util
    import java.util.ArrayList; // importeert enkel de ArrayList klasse
\end{codeblock}

\subsection{Typische Java klasse structuur \& waarden}
Een typische Java klasse heeft de volgende structuur:
\begin{codeblock}[Hoe java doc starten?]{java}
    import java.io.*; //zoals File, FileReader, BufferedReader, etc

    import java.util.*; //zoals ArrayList, Scanner, Random, etc


    public class KlasseNaam  // klasse definitie
        {
            // opstellen van variabelen (attributen)
            private int attribuut1;
            private String attribuut2;
            private float attribuut3; // niet alle attributen moeten in de constructor zitten
            public int attribuut4; // kan ook public zijn

            // constructor (Je creeert een instance van een klasse met een constructor)
            public KlasseNaam(int attribuut1, String attribuut2)
            {
                    this.attribuut1 = attribuut1;
                    this.attribuut2 = attribuut2;
                }

            // methoden
            public void methode1()
            {
                    // methode implementatie
                }

            public int methode2(String parameter)
            {
                    // methode implementatie
                    return 0;
                }

            // Belangrijk: na elke instructie een puntkomma ;
        }
\end{codeblock}

\subsubsection{Datatypes}
Een variabele heeft altijd een type. Dat type bepaalt wat er in mag en wat je
ermee mag doen.

\begin{center}
    \begin{tabular}{@{}llp{6.2cm}@{}}
        \toprule
        \textbf{Type}               & \textbf{Voorbeeld}         & \textbf{Wanneer}                                    \\
        \midrule
        \texttt{int}                & \texttt{5}                 & gehele getallen: aantal, index, jaartal             \\
        \texttt{long}               & \texttt{5L}                & heel grote gehele getallen                          \\
        \texttt{double}             & \texttt{5.0}               & kommagetallen (standaardkeuze)                      \\
        \texttt{float}              & \texttt{5.0f}              & kommagetallen, minder precies; let op de \texttt{f} \\
        \texttt{boolean}            & \texttt{true}              & ja/nee, resultaat van een voorwaarde                \\
        \texttt{char}               & \texttt{'a'}               & \'e\'en teken, \textbf{enkele} aanhalingstekens     \\
        \texttt{String}             & \texttt{"abc"}             & tekst, \textbf{dubbele} aanhalingstekens            \\
        \midrule
        \texttt{int[]}              & \texttt{new int[5]}        & rij getallen met vaste lengte                       \\
        \texttt{int[][]}            & \texttt{new int[3][4]}     & tabel (rijen en kolommen)                           \\
        \texttt{ArrayList<Integer>} & \texttt{new ArrayList<>()} & rij die kan groeien en krimpen                      \\
        \bottomrule
    \end{tabular}
\end{center}

\begin{warningbox}[Twee klassiekers]
    \texttt{5 / 2} geeft \texttt{2}, niet \texttt{2.5}. Deel je twee \texttt{int}'s,
    dan is het resultaat weer een \texttt{int} en valt alles na de komma weg. Wil je
    een kommagetal, maak dan \'e\'en van de twee decimaal:
    \texttt{5 / 2.0} of \texttt{(double) a / b}.

    \medskip
    \texttt{'F'} is een \texttt{char}, \texttt{"F"} is een \texttt{String}. Die twee
    zijn niet uitwisselbaar.
\end{warningbox}

\textbf{Public} betekent dat andere klassen deze variabele of methode
kunnen gebruiken.
\textbf{Private} betekent dat alleen deze klasse deze variabele of methode
kan gebruiken.

\subsection{Methoden}

een methode zit in je klasse. Hiermee
kun je functionaliteit toevoegen aan je klasse.

Een methodekop bestaat altijd uit vier stukken, in deze volgorde:
\keyterm{zichtbaarheid} --- \keyterm{returntype} --- \keyterm{naam} ---
\keyterm{parameters}.

\begin{codeblock}[methode voorbeeld]{java}
    // public   int      methodeOptelling(int a, int b)
    // zichtb.  return   naam            parameters
    public int methodeOptelling(int parameter1, int parameter2)
    {
            int resultaat = parameter1 + parameter2;
            return resultaat; // geeft het resultaat terug
        }

    // een methode met returntype 'void' geeft niets terug:
    public void methodePrint(String tekst)
    {
            System.out.println(tekst); // print de tekst naar de console
        }
\end{codeblock}

Je mag elk type teruggeven: een \texttt{int}, een \texttt{boolean}, een
\texttt{String}, maar ook een zelfgemaakte klasse (\texttt{Persoon}) of een hele
collectie (\texttt{ArrayList<Persoon>}).

\begin{warningbox}[Elk pad moet iets teruggeven]
    Heeft je methode een returntype (dus niet \texttt{void}), dan moet \'elke
    mogelijke weg door je code op een \texttt{return} eindigen. Ook de \texttt{else},
    ook de \texttt{catch}. Anders krijg je de compileerfout
    \textit{missing return statement}. Een veelgebruikte truc is onderaan een
    foutwaarde teruggeven: \texttt{return null;} voor een object,
    \texttt{return -1;} voor een aantal.
\end{warningbox}

\section{Strings vergelijken en bewerken}
\label{sec:strings}

\subsection{== of .equals()?}

Een variabele van het type \texttt{int}, \texttt{char}, \texttt{double} of
\texttt{boolean} bevat de \textbf{waarde zelf}. Een variabele van het type
\texttt{String}, of van eender welke klasse, bevat de waarde \textbf{niet}: ze
bevat het \keyterm{adres} van het object dat ergens anders in het geheugen staat.
Je kan zo'n variabele zien als een pijl naar een doos.

\begin{center}
    \begin{tikzpicture}[font=\small\ttfamily, node distance=1cm]
        \node[draw, thick, minimum width=1.6cm, minimum height=0.7cm] (a) at (0,1.4) {a};
        \node[draw, thick, minimum width=1.6cm, minimum height=0.7cm] (b) at (0,0) {b};
        \node[draw, thick, fill=schoolBlue!12, minimum width=1.8cm, minimum height=0.7cm] (o1) at (4.2,1.4) {\textquotedbl Bob\textquotedbl};
        \node[draw, thick, fill=schoolBlue!12, minimum width=1.8cm, minimum height=0.7cm] (o2) at (4.2,0) {\textquotedbl Bob\textquotedbl};
        \draw[-{Latex[length=2.5mm]}, thick] (a) -- (o1);
        \draw[-{Latex[length=2.5mm]}, thick] (b) -- (o2);
        \node[font=\scriptsize\sffamily, anchor=west] at (5.6,1.4) {twee \emph{aparte} objecten};
        \node[font=\scriptsize\sffamily, anchor=west] at (5.6,0)   {met dezelfde inhoud};
        \node[font=\small\sffamily, anchor=west] at (0.1,-1.1)
        {\texttt{a == b} is \textbf{false} (andere adressen), \texttt{a.equals(b)} is \textbf{true} (zelfde tekst)};
    \end{tikzpicture}
\end{center}

\begin{itemize}
    \item \texttt{==} vergelijkt de \textbf{pijlen}: wijzen de twee variabelen
          naar hetzelfde object?
    \item \texttt{.equals()} vergelijkt de \textbf{inhoud}: staat er dezelfde
          tekst in?
\end{itemize}

Je wil bijna altijd de inhoud weten, dus bijna altijd \texttt{.equals()}.

\begin{codeblock}[== of equals]{java}
    String a = new String("Bob");
    String b = new String("Bob");

    if (a == b)      { }   // FALSE: twee verschillende objecten
    if (a.equals(b)) { }   // TRUE : zelfde tekst

    int x = 5;
    int y = 5;
    if (x == y)      { }   // TRUE : bij int is == wel juist
\end{codeblock}


\begin{center}
    \begin{tabular}{@{}llp{6.4cm}@{}}
        \toprule
        \textbf{Type}                                 & \textbf{Vergelijk met}         & \textbf{Opmerking}                        \\
        \midrule
        \texttt{int}, \texttt{char}, \texttt{boolean} & \texttt{==}                    & bevat de waarde zelf                      \\
        \texttt{double}, \texttt{float}               & \texttt{Math.abs(a-b) < 0.001} & afrondingsfouten, zie hieronder           \\
        \texttt{String}                               & \texttt{.equals()}             & of \texttt{.equalsIgnoreCase()}           \\
        eigen klasse                                  & een veld vergelijken           & bv.\ \texttt{a.getId().equals(b.getId())} \\
        \texttt{null}                                 & \texttt{==}                    & \texttt{if (p == null)} is w\'el juist    \\
        \bottomrule
    \end{tabular}
\end{center}

\begin{warningbox}[Twee vervolgvallen]
    \textbf{Kommagetallen.} \texttt{0.1 + 0.2 == 0.3} is \texttt{false}. Een
    \texttt{double} kan niet elk getal exact voorstellen. Vergelijk daarom op
    ``dicht genoeg bij elkaar'': \texttt{if (Math.abs(a - b) < 0.001)}.

    \medskip
    \textbf{Volgorde bij null.} \texttt{p.getNaam().equals("Bob")} crasht met een
    \textit{NullPointerException} als \texttt{p} zelf \texttt{null} is, bijvoorbeeld
    omdat je zoekmethode niets vond. Check dus eerst:
    \texttt{if (p != null \&\& p.getNaam().equals("Bob"))}. Java stopt bij
    \texttt{\&\&} met evalueren zodra het linkse deel \texttt{false} is, dus die
    volgorde is veilig.
\end{warningbox}

\subsection{Eigen objecten vergelijken}

Twee objecten van je eigen klasse vergelijken met \texttt{==} werkt alleen als
het letterlijk hetzelfde object is. Wil je weten of het ``dezelfde persoon'' is,
vergelijk dan het veld dat uniek is:

\begin{codeblock}[eigen objecten vergelijken]{java}
    // FOUT: vergelijkt adressen
    if (kandidaat1 == kandidaat2) { }

    // JUIST: vergelijk het kenmerk dat uniek is
    if (kandidaat1.getMobielNummer().equals(kandidaat2.getMobielNummer())) { }
\end{codeblock}

\texttt{lijst.contains(x)} en \texttt{lijst.remove(x)} roepen \texttt{equals()}
op. Schrijf jij zelf geen \texttt{equals()} in je klasse, dan gebruikt Java de
standaardversie, en die vergelijkt adressen. \texttt{contains()} vindt dan alleen
exact hetzelfde object terug, niet een object met dezelfde gegevens. Schrijf in
dat geval je eigen zoekmethode met een lus.

\subsection{Nuttige String-methodes}

\begin{center}
    \begin{tabular}{@{}llp{5.9cm}@{}}
        \toprule
        \textbf{Methode}             & \textbf{Geeft terug} & \textbf{Doet}                                             \\
        \midrule
        \texttt{length()}            & \texttt{int}         & aantal tekens                                             \\
        \texttt{charAt(i)}           & \texttt{char}        & teken op index \texttt{i}                                 \\
        \texttt{substring(a, b)}     & \texttt{String}      & van index \texttt{a} tot \textbf{net v\'o\'or} \texttt{b} \\
        \texttt{substring(a)}        & \texttt{String}      & vanaf index \texttt{a} tot het einde                      \\
        \texttt{indexOf("x")}        & \texttt{int}         & positie, of \texttt{-1} als het er niet in zit            \\
        \texttt{contains("x")}       & \texttt{boolean}     & komt het erin voor?                                       \\
        \texttt{startsWith("04")}    & \texttt{boolean}     & begint ermee?                                             \\
        \texttt{equals(t)}           & \texttt{boolean}     & zelfde tekst?                                             \\
        \texttt{equalsIgnoreCase(t)} & \texttt{boolean}     & zelfde tekst, hoofdletters negeren                        \\
        \texttt{toLowerCase()}       & \texttt{String}      & alles klein                                               \\
        \texttt{trim()}              & \texttt{String}      & spaties vooraan en achteraan weg                          \\
        \texttt{isEmpty()}           & \texttt{boolean}     & is de tekst leeg?                                         \\
        \texttt{split(";")}          & \texttt{String[]}    & opsplitsen op een scheidingsteken                         \\
        \bottomrule
    \end{tabular}
\end{center}

\begin{warningbox}[Een String verandert nooit]
    \texttt{toLowerCase()}, \texttt{trim()} en \texttt{substring()} passen de
    originele String \textbf{niet} aan: ze geven een \textbf{nieuwe} String terug. Je
    moet het resultaat dus opvangen.

    \begin{codeblock}[fout en juist]{java}
        naam.toLowerCase();          // doet niets zichtbaar
        naam = naam.toLowerCase();   // zo wel
    \end{codeblock}
\end{warningbox}

\begin{codeblock}[String in de praktijk]{java}
    String nr = "0470123456";

    nr.length()            // 10
    nr.charAt(0)           // '0'  (een char, geen String)
    nr.substring(0, 2)     // "04" : index 0 en 1, NIET index 2
    nr.startsWith("04")    // true

    // een regel uit een bestand opsplitsen
    String regel = "Bob;1999;M";
    String[] delen = regel.split(";");
    String naam = delen[0];                    // "Bob"
    int jaar = Integer.parseInt(delen[1]);     // 1999
    char geslacht = delen[2].charAt(0);        // 'M'
\end{codeblock}

\subsection{Omzetten tussen tekst en getallen}

\begin{codeblock}[omzetten]{java}
    int i = Integer.parseInt("42");        // String -> int
    double d = Double.parseDouble("4.2");  // String -> double
    String s = String.valueOf(42);         // int -> String
    String t = "" + 42;                    // kortere truc

    // losse tekens onderzoeken
    Character.isDigit('7')    // true
    Character.isLetter('a')   // true
\end{codeblock}

\begin{warningbox}[Plus betekent twee dingen]
    Bij getallen telt \texttt{+} op, bij tekst plakt \texttt{+} aan elkaar. Java
    werkt van links naar rechts, dus:

    \medskip
    \texttt{1 + 2 + "a"} geeft \texttt{"3a"} (eerst optellen, dan plakken) \\
    \texttt{"a" + 1 + 2} geeft \texttt{"a12"} (vanaf de tekst wordt alles plakken)

    \medskip
    Zet haakjes als je twijfelt: \texttt{"totaal: " + (a + b)}.
\end{warningbox}

\begin{examenbox}
    Onthoud de drie vormen. Bij een array is het \texttt{length} zonder haakjes, bij
    een String \texttt{length()} met haakjes, en bij een lijst een heel ander woord:
    \texttt{size()}. Verwarring hiertussen is een klassieke compileerfout.
\end{examenbox}

\section{If, for en while}
\subsection{If statements}
\begin{codeblock}[if statement]{java}
    if (voorwaarde) {
            // code die uitgevoerd wordt als de voorwaarde waar is
        } else if (andereVoorwaarde)
    {
            // code die uitgevoerd wordt als de andere voorwaarde waar is
        } else
        {
            // code die uitgevoerd wordt als geen van de voorwaarden waar is
        }

    if en else zijn een sort scan, dus eerst wordt if gedaan en als
    dat niet waar is, gaat die naar else if en zo verder.
    Moest je nu gewoon if statements gebruiken, dan worden ze allemaal
    uitgevoerd.

    Dit zijn belangrijke voorwaarden:
    ==  // gelijk aan
    !=  // niet gelijk aan
    >   // groter dan
    <   // kleiner dan
    >=  // groter dan of gelijk aan
    <=  // kleiner dan of gelijk aan
    &&  // en
    ||  // of
\end{codeblock}

\subsection{For loops}
\begin{codeblock}[for loop]{java}
    for (init; voorwaarde; increment)
    {
            // code die herhaald wordt
        }

    // voorbeeld: print getallen van 0 tot 9
    for (int i = 0; i < 10; i++) // i begint bij 0, zolang i kleiner is dan 10, verhoog i met 1
    //vergeet de ; niet tussen de statements (init, voorwaarde, increment)
    {
            System.out.println(i); //output is 0 1 2 3 4 5 6 7 8 9
        }
\end{codeblock}

\subsection{While loops}
Een while-lus draait zolang de voorwaarde waar is.

\begin{codeblock}[while loop]{java}
    int i = 0; // initialisatie
    while (i < voorwaarde) // zolang i kleiner is dan de voorwaarde
        {
            // code die herhaald wordt zolang de voorwaarde waar is
            i++; // vergeet niet i aan te passen om een oneindige loop te vermijden
            i--; // of i verlagen
        }
\end{codeblock}


\begin{conceptbox}[Welke lus kies je?]
    \textbf{for-each} als je gewoon elk element \'e\'en keer wil bekijken. Korter en
    je kan je niet vergissen in de grenzen.

    \medskip
    \textbf{gewone for} als je de \textbf{index} nodig hebt: om op positie
    \texttt{i} te schrijven, om twee arrays tegelijk te doorlopen, of om achterwaarts
    te gaan.

    \medskip
    \textbf{while} als je op voorhand niet weet hoeveel rondes je nodig hebt (lezen
    tot het bestand op is, blijven proberen tot iets uniek is).
\end{conceptbox}

\subsection{Arrays}

een array is wiskundig een vector.
Een array van strings zijn gewoon meerdere teksten in 1 variabele.

\begin{codeblock}[array]{java}
    // array van gehele getallen met 5 elementen
    int[] getallen = new int[5];

    // array van strings met 3 elementen
    String[] namen = new String[3]; // een array zijn lengte is vast. pas op want dit is niet de index.

    // array initialisatie met waarden
    int[] cijfers = {90, 85, 78, 92, 88};

    // toegang tot array elementen
    int eersteGetal = getallen[0]; // eerste element (index 0)
    String eersteNaam = namen[0]; // eerste element (index 0)

    //belangrijke array functies
    int lengte = cijfers.length; // lengte van de array
    int laatsteIndex = cijfers.length - 1; // laatste index van de array
    int middelsteIndex = cijfers.length / 2; // middelste index van de array

    // itereren over een array
    for (int i = 0; i < cijfers.length; i++)
    {
            System.out.println(cijfers[i]);
        }

    // De for-each loop. Deze moet je zeker kennen voor ArrayLists,
    // maar is ook bij arrays handig.
    // Lees ze als: "voor elk cijfer in cijfers"
    for (int cijfer : cijfers) // let op de : , geen ;
    {
    System.out.println(cijfer);
    }
\end{codeblock}

\begin{warningbox}[Index loopt van 0 tot length-1]
    Een array van lengte 5 heeft indexen 0, 1, 2, 3, 4. Dus \texttt{i < length} en
    nooit \texttt{i <= length}. Doe je dat toch, dan crasht je programma met
    \textit{ArrayIndexOutOfBoundsException}.


    \medskip
    \begin{tabular}{@{}lll@{}}
        \toprule
        \textbf{Doel}       & \textbf{Juist}                  & \textbf{Fout}               \\
        \midrule
        door een array      & \texttt{i < a.length}           & \texttt{<=} geeft een crash \\
        door een String     & \texttt{i < s.length()}         & \texttt{<=} geeft een crash \\
        8 tekens verzamelen & \texttt{while (s.length() < 8)} & \texttt{<= 8} geeft er 9    \\
        \bottomrule
    \end{tabular}
\end{warningbox}


\subsection{Multi dim arrays}

\begin{codeblock}[multi dim array]{java}
    // 2D array (matrix) van gehele getallen met 3 rijen en 4 kolommen
    int[][] matrix = new int[3][4];

    // initialisatie van een 2D array
    int[][] tabel = {
    {1, 2, 3, 4},
    {5, 6, 7, 8},
    {9, 10, 11, 12}
    };

    // toegang tot elementen in een 2D array
    int element = tabel[1][2]; // element in de tweede rij en derde kolom (waarde is 7)

    // itereren met indexen
    for (int i = 0; i < tabel.length; i++)      // .length = aantal rijen
        {
            for (int j = 0; j < tabel[i].length; j++) // lengte van die ene rij
                {
                    System.out.print(tabel[i][j] + " ");
                }
            System.out.println(); // nieuwe regel na elke rij
        }
    // output: 1 2 3 4 / 5 6 7 8 / 9 10 11 12

    // itereren met for-each
    for (int[] rij : tabel)   // een element van een int[][] is een int[] !
    {
    for (int getal : rij) // pas hier krijg je losse getallen
        {
            System.out.print(getal + " ");
        }
    }
\end{codeblock}

\begin{warningbox}[Veelgemaakte fout bij 2D en for-each]
    \texttt{for (int i : tabel)} compileert \textbf{niet}. Een 2D-array is een array
    \emph{van arrays}: elk element is zelf een \texttt{int[]}, geen \texttt{int}. Je
    hebt dus twee niveaus nodig: eerst de rij, dan het getal.

    \medskip
    Verwar ook \texttt{tabel.length} (aantal rijen) niet met
    \texttt{tabel[i].length} (aantal kolommen in rij \texttt{i}).
\end{warningbox}

\subsection{ArrayLists}

Een array heeft een \textbf{vaste} lengte: eens \texttt{new int[5]} kan er nooit
een zesde bij. Een \texttt{ArrayList} groeit en krimpt vanzelf. Daarom gebruik je
in dit vak bijna altijd een \texttt{ArrayList}: je weet op voorhand niet hoeveel
kandidaten, tickets of mieren er zullen zijn.

\begin{center}
    \begin{tabular}{@{}llp{5.6cm}@{}}
        \toprule
        \textbf{Methode}     & \textbf{Geeft terug} & \textbf{Doet}                  \\
        \midrule
        \texttt{add(e)}      & \texttt{boolean}     & achteraan toevoegen            \\
        \texttt{add(i, e)}   & \texttt{void}        & invoegen op index \texttt{i}   \\
        \texttt{get(i)}      & element              & element op index \texttt{i}    \\
        \texttt{set(i, e)}   & oud element          & vervangen op index \texttt{i}  \\
        \texttt{remove(i)}   & verwijderd elem.     & verwijderen op \textbf{index}  \\
        \texttt{remove(e)}   & \texttt{boolean}     & verwijderen op \textbf{object} \\
        \texttt{size()}      & \texttt{int}         & aantal elementen               \\
        \texttt{isEmpty()}   & \texttt{boolean}     & is de lijst leeg?              \\
        \texttt{contains(e)} & \texttt{boolean}     & zit \texttt{e} erin?           \\
        \texttt{indexOf(e)}  & \texttt{int}         & index of \texttt{-1}           \\
        \texttt{clear()}     & \texttt{void}        & alles wissen                   \\
        \bottomrule
    \end{tabular}
\end{center}

\begin{warningbox}[remove(i) of remove(e)?]
    Bij een \texttt{ArrayList<Integer>} is \texttt{remove(2)} het element op
    \textbf{index 2}, maar \texttt{remove(Integer.valueOf(2))} is het \textbf{getal
        2}. Java kiest op basis van het type, niet van wat jij bedoelt.
\end{warningbox}

\begin{codeblock}[arraylist]{java}
    import java.util.ArrayList; //importeer de ArrayList klasse

    // maak een ArrayList van Strings

    private ArrayList<String> namen;

    in constructor:
    public MijnKlasse()
    {
            namen = new ArrayList<String>();
        }

    // voeg elementen toe aan de ArrayList
    namen.add("Alice");
    namen.add("Bob");
    namen.add("Charlie");

    // toegang tot elementen in de ArrayList
    String eersteNaam = namen.get(0); // eerste element (index 0)

    // verwijder een element uit de ArrayList
    namen.remove(1); // verwijdert het element op index 1 ("Bob")
    namen.remove("Bob"); // verwijdert het element "Bob" als het bestaat

    // grootte van de ArrayList
    int grootte = namen.size(); // aantal elementen in de ArrayList

    // itereren over een ArrayList
    for (String naam : namen)
    {
            System.out.println(naam);
            //output: Alice Charlie
        }

    //arraylists kunnen ook instancies bevatten van
    andere klassen,
    bijvoorbeeld een arraylist van personen:

    private ArrayList<Persoon> personen;

    in constructor:
    public MijnKlasse()
    {
            personen = new ArrayList<Persoon>();
        }

    // voeg een persoon toe aan de ArrayList
    public void voegPersoonToe(String naam, int leeftijd)
    {
            personen.add(new Persoon(naam, leeftijd)); //persoon(constructor)
        }

    // Je kunt in de lus de methodes van de Persoon-klasse gebruiken.
\end{codeblock}

\subsubsection{Verwijderen terwijl je door de lijst gaat}

Onderstaande methode \textbf{ziet er
    logisch uit maar werkt niet}:

\begin{codeblock}[FOUT: verwijderen in een for-each]{java}
    public ArrayList<Persoon> verwijderOuderePersonen(int grens) //methode die een ArrayList teruggeeft
        {
            ArrayList<Persoon> verwijderd = new ArrayList<Persoon>();
            for (Persoon persoon : personen)
            {
                    if (persoon.getLeeftijd() > grens)
                    {
                            verwijderd.add(persoon);
                            personen.remove(persoon); // <-- crasht
                        }
                }
            // <-- en hier ontbreekt bovendien de return
        }
\end{codeblock}

Twee fouten:
\begin{enumerate}
    \item De methode belooft een \texttt{ArrayList<Persoon>} terug te geven maar
          doet dat nergens. Dat compileert niet:
          \textit{missing return statement}.
    \item Je verandert de lijst terwijl de for-each er nog over loopt. Wat er dan
          gebeurt, staat in de figuur hieronder:
          \textit{ConcurrentModificationException}.
\end{enumerate}

\begin{center}
\begin{tikzpicture}[
  vak/.style={draw=black!55, minimum width=13mm, minimum height=6.5mm,
              font=\ttfamily\footnotesize, inner sep=0pt},
  x=13mm, y=-17mm]

\definecolor{wegKleur}{RGB}{255,214,214}
\definecolor{misKleur}{RGB}{255,226,170}

\node[anchor=east, font=\footnotesize\bfseries] at (-0.15, 0) {stap 1};
\foreach \t [count=\x from 0] in {Ann, Bob, Cas, Dora}
    \node[vak, fill=white] at (\x+0.5, 0) {\t};
\draw[-{Stealth}, schoolBlue, thick] (1.5, 0.42) -- (1.5, 0.16);
\node[font=\tiny, schoolBlue, anchor=north] at (1.5, 0.44) {for-each staat op plaats 1};

\node[anchor=east, font=\footnotesize\bfseries] at (-0.15, 1) {stap 2};
\foreach \t [count=\x from 0] in {Ann, Cas, Dora}
    \node[vak, fill=white] at (\x+0.5, 1) {\t};
\node[vak, fill=wegKleur, dashed] at (3.5, 1) {};
\node[font=\tiny, anchor=west] at (4.2, 1) {Bob is weg, alles schuift een plaats op};
\draw[-{Stealth}, schoolBlue, thick] (1.5, 1.42) -- (1.5, 1.16);
\node[font=\tiny, schoolBlue, anchor=north] at (1.5, 1.44) {de pijl staat nog altijd op plaats 1};

\node[anchor=east, font=\footnotesize\bfseries] at (-0.15, 2) {stap 3};
\node[vak, fill=white]    at (0.5, 2) {Ann};
\node[vak, fill=misKleur] at (1.5, 2) {Cas};
\node[vak, fill=white]    at (2.5, 2) {Dora};
\draw[-{Stealth}, schoolBlue, thick] (2.5, 2.42) -- (2.5, 2.16);
\node[font=\tiny, schoolBlue, anchor=north] at (2.6, 2.44) {volgende ronde: plaats 2};
\node[font=\tiny, anchor=west] at (4.2, 2) {Cas wordt overgeslagen $\rightarrow$ Java gooit};
\node[font=\tiny\itshape, anchor=west] at (4.2, 2.22) {ConcurrentModificationException};

\end{tikzpicture}
\end{center}

De oplossing is een \keyterm{Iterator}. Dat is het object dat de
for-each achter de schermen ook gebruikt, maar nu heb jij het zelf in handen. Zijn
\texttt{remove()} verwijdert het laatst opgehaalde element \'en houdt de positie
correct.

\begin{codeblock}[JUIST: verwijderen met een Iterator]{java}
    import java.util.Iterator;

    public ArrayList<Persoon> verwijderOuderePersonen(int grens)
    {
    ArrayList<Persoon> verwijderd = new ArrayList<Persoon>();

    Iterator<Persoon> it = personen.iterator();
    while (it.hasNext())        // is er nog een volgend element?
    {
    Persoon persoon = it.next();   // haal het op EN ga een stap verder
    if (persoon.getLeeftijd() > grens)
    {
            verwijderd.add(persoon);
            it.remove();        // via de iterator, niet via personen
        }
    }
    return verwijderd;
    }
\end{codeblock}

\begin{warningbox}[it.next() maar een keer per ronde]
    \texttt{it.next()} mag je \textbf{maar \'e\'en keer per ronde} oproepen: elke
    oproep schuift een plaats op. Wil je de waarde twee keer gebruiken, zet ze dan
    eerst in een variabele, zoals hierboven.
\end{warningbox}

\begin{conceptbox}[Wanneer heb je een Iterator nodig?]
    Alleen als je \textbf{verwijdert tijdens het doorlopen}. Wil je enkel lezen,
    tellen of optellen, gebruik dan gewoon de for-each: die is korter en duidelijker.
\end{conceptbox}

\subsubsection{Verwijderen zonder Iterator: de removelijst}

Je gaat eerste een removellijst maken afhankelijk van je conditie en je gaat dan met een
for loop itereren en elk element in de originele lijst verwijderen.

\begin{codeblock}[JUIST: verwijderen via een hulplijst]{java}
    public int verwijderOuderePersonen(int grens)
    {
            ArrayList<Persoon> weg = new ArrayList<Persoon>();

            // ronde 1: enkel lezen uit personen, niets aanraken
            for (Persoon p : personen)
            {
                    if (p.getLeeftijd() > grens)
                    {
                            weg.add(p);
                        }
                }

            // ronde 2: nu loop je over weg en verander je personen
            for (Persoon p : weg)
            {
                    personen.remove(p);
                }

            return weg.size();   // de teller krijg je gratis
        }
\end{codeblock}

Waarom mag dit w\'el? In ronde 1 lees je alleen. In ronde 2 loop je over
\texttt{weg} en pas je \texttt{personen} aan --- twee \textbf{verschillende}
lijsten. De for-each houdt alleen zijn plaats bij in de lijst waar hij over
loopt, en die verandert niet.

De laatste lus kan ook in \'e\'en regel:

\begin{codeblock}[Korter]{java}
    personen.removeAll(weg);
\end{codeblock}

\begin{center}
    \begin{tabular}{@{}lll@{}}
        \toprule
                       & \textbf{Iterator}           & \textbf{Hulplijst}                \\
        \midrule
        Aantal lussen  & 1                           & 2                                 \\
        Extra geheugen & nee                         & ja, een tweede lijst              \\
        Teller         & zelf \texttt{n++}           & \texttt{weg.size()}               \\
        Extra import   & \texttt{java.util.Iterator} & geen                              \\
        Verwijdert     & exact dit element           & het \emph{eerste} gelijke element \\
        \bottomrule
    \end{tabular}
\end{center}

\begin{warningbox}[Klein verschil dat je moet kennen]
    \texttt{personen.remove(p)} zoekt het \textbf{eerste} element dat gelijk is aan
    \texttt{p}. Zolang je met dezelfde objecten werkt, is dat precies het juiste.
    Heeft je klasse een eigen \texttt{equals()}, dan kan hij een ander object pakken
    dat toevallig gelijk is. De \texttt{Iterator} heeft dat probleem nooit: die
    verwijdert exact het element waar hij staat.
\end{warningbox}


\subsection{Sorteren}

Java sorteert een \texttt{ArrayList} van je eigen objecten niet vanzelf: het weet
niet welk veld belangrijk is. Je schrijft de sortering dus zelf, met een
selectiesortering.

\medskip
De buitenste lus (\texttt{i}) zegt welke plaats je invult, de binnenste lus
(\texttt{j}) zoekt wie daar hoort.

\begin{center}
\begin{tikzpicture}[
  vak/.style={draw=black!55, minimum width=13mm, minimum height=6.5mm,
              font=\ttfamily\footnotesize, inner sep=0pt},
  x=13mm, y=-15mm]

\definecolor{sortVast}{RGB}{215,232,244}
\definecolor{sortKlein}{RGB}{255,226,170}
\definecolor{sortKlaar}{RGB}{223,240,225}

% ronde 1
\node[anchor=east, font=\footnotesize\bfseries] at (-0.15, 0) {ronde 1};
\node[vak, fill=sortVast]  at (0.5, 0) {delta};
\node[vak, fill=white]     at (1.5, 0) {bravo};
\node[vak, fill=white]     at (2.5, 0) {echo};
\node[vak, fill=sortKlein] at (3.5, 0) {alfa};
\node[font=\tiny, anchor=south] at (0.5, -0.20) {i=0};
\draw[-{Stealth}, schoolBlue] (1.2, -0.20) -- (3.8, -0.20);
\node[font=\tiny, schoolBlue, anchor=south] at (2.5, -0.26) {j};
\draw[{Stealth}-{Stealth}, red!65!black]
      (0.5, 0.24) to[bend right=20] (3.5, 0.24);
\node[font=\tiny, anchor=west] at (4.15, 0) {wissel plaats 0 en 3};

% ronde 2
\node[anchor=east, font=\footnotesize\bfseries] at (-0.15, 1) {ronde 2};
\node[vak, fill=sortKlaar] at (0.5, 1) {alfa};
\node[vak, fill=sortVast]  at (1.5, 1) {bravo};
\node[vak, fill=white]     at (2.5, 1) {echo};
\node[vak, fill=white]     at (3.5, 1) {delta};
\node[font=\tiny, anchor=south] at (1.5, 0.80) {i=1};
\draw[-{Stealth}, schoolBlue] (2.2, 0.80) -- (3.8, 0.80);
\node[font=\tiny, schoolBlue, anchor=south] at (3.0, 0.74) {j};
\node[font=\tiny, anchor=west] at (4.15, 1) {bravo is al de kleinste, niets wisselen};

% ronde 3
\node[anchor=east, font=\footnotesize\bfseries] at (-0.15, 2) {ronde 3};
\node[vak, fill=sortKlaar] at (0.5, 2) {alfa};
\node[vak, fill=sortKlaar] at (1.5, 2) {bravo};
\node[vak, fill=sortVast]  at (2.5, 2) {echo};
\node[vak, fill=sortKlein] at (3.5, 2) {delta};
\node[font=\tiny, anchor=south] at (2.5, 1.80) {i=2};
\draw[-{Stealth}, schoolBlue] (3.2, 1.80) -- (3.8, 1.80);
\node[font=\tiny, schoolBlue, anchor=south] at (3.5, 1.74) {j};
\draw[{Stealth}-{Stealth}, red!65!black]
      (2.5, 2.24) to[bend right=26] (3.5, 2.24);
\node[font=\tiny, anchor=west] at (4.15, 2) {wissel plaats 2 en 3};

% klaar
\node[anchor=east, font=\footnotesize\bfseries] at (-0.15, 3) {klaar};
\node[vak, fill=sortKlaar] at (0.5, 3) {alfa};
\node[vak, fill=sortKlaar] at (1.5, 3) {bravo};
\node[vak, fill=sortKlaar] at (2.5, 3) {delta};
\node[vak, fill=sortKlaar] at (3.5, 3) {echo};

\node[font=\tiny, anchor=west] at (-0.1, 3.6)
  {\tikz\node[draw=black!55, fill=sortVast, minimum width=3.5mm, minimum height=2.5mm, inner sep=0pt]{};\ plaats i die je invult
   \quad
   \tikz\node[draw=black!55, fill=sortKlein, minimum width=3.5mm, minimum height=2.5mm, inner sep=0pt]{};\ kleinste gevonden
   \quad
   \tikz\node[draw=black!55, fill=sortKlaar, minimum width=3.5mm, minimum height=2.5mm, inner sep=0pt]{};\ staat vast};

\end{tikzpicture}
\end{center}

\subsubsection{Twee objecten vergelijken}

Een \texttt{Item} kan je niet rechtstreeks vergelijken. De \texttt{String} erin
wel. Haal die met de getter op en gebruik \texttt{compareTo}:

\begin{center}
\begin{tabular}{@{}ll@{}}
\toprule
\texttt{a.compareTo(b)} & \textbf{Betekenis} \\
\midrule
negatief & \texttt{a} komt alfabetisch eerst \\
0        & de twee zijn gelijk \\
positief & \texttt{b} komt eerst \\
\bottomrule
\end{tabular}
\end{center}

\begin{codeblock}[alfabetisch sorteren]{java}
    public void sorteerAlfabetisch(ArrayList<Item> lijst)
    {
        for (int i = 0; i < lijst.size(); i++)
        {
            int kleinste = i;              // voorlopig de kleinste

            for (int j = i + 1; j < lijst.size(); j++)
            {
                String a = lijst.get(j).getCode();
                String b = lijst.get(kleinste).getCode();
                if (a.compareTo(b) < 0)
                {
                    kleinste = j;
                }
            }

            Item temp = lijst.get(i);      // de twee plaatsen omwisselen
            lijst.set(i, lijst.get(kleinste));
            lijst.set(kleinste, temp);
        }
    }
\end{codeblock}

De \texttt{temp}-variabele is nodig. Schrijf je meteen
\texttt{lijst.set(i, lijst.get(kleinste))}, dan is wat op plaats \texttt{i} stond
weg voor je het kon bewaren.

\begin{warningbox}[Sorteren kan niet met een for-each]
Sorteren betekent \texttt{lijst.set(i, \dots)}: je schrijft op een \textbf{plaats}
terug. Een for-each geeft je alleen de waarde, niet het plaatsnummer. En dit doet
niets aan de lijst:
\begin{center}
\texttt{for (Item i : lijst) \{ i = anderItem; \}}
\end{center}
\texttt{i} is een kopie van de verwijzing. Je verzet die kopie, de lijst blijft
zoals ze was. Bij sorteren gebruik je dus altijd een gewone for.
\end{warningbox}

\subsubsection{Sorteren op twee dingen tegelijk}

Vraag uit een examen: sorteer zo dat alle objecten met \texttt{true} vooraan
staan, en binnen elke groep alfabetisch. Splits de lijst, sorteer elk deel apart,
plak ze weer aan elkaar.

\begin{codeblock}[splitsen, sorteren, samenvoegen]{java}
    public ArrayList<Item> gesorteerd()
    {
        ArrayList<Item> waar = new ArrayList<Item>();
        ArrayList<Item> nietWaar = new ArrayList<Item>();

        // 1. splitsen (hier lees je alleen, dus for-each mag)
        for (Item i : items)
        {
            if (i.isActief())
            {
                waar.add(i);
            }
            else
            {
                nietWaar.add(i);
            }
        }

        // 2. elk deel apart sorteren
        sorteerAlfabetisch(waar);
        sorteerAlfabetisch(nietWaar);

        // 3. samenvoegen: eerst de true's, dan de false's
        ArrayList<Item> resultaat = new ArrayList<Item>();
        resultaat.addAll(waar);
        resultaat.addAll(nietWaar);
        return resultaat;
    }
\end{codeblock}

\texttt{addAll} plakt een hele lijst achteraan een andere. Dat spaart je een lus
uit.

\begin{examenbox}
Splitsen en weer samenvoegen is makkelijker te schrijven dan een sortering met
twee voorwaarden in \'e\'en lus. Zeker onder tijdsdruk: elke stap is apart te
testen.
\end{examenbox}


\section{Class extends}
Met \texttt{extends} bouw je een klasse voort op een andere. Je schrijft
\texttt{public class Hond extends Dier}. \texttt{Hond} krijgt dan alle velden en
methodes van \texttt{Dier} erbij, zonder dat je ze overtypt. In het Engels heet
dat \textit{inheritance}.

\begin{codeblock}[class extends]{java}
    // Bestaande klasse
    public class Dier
        {

            private String naam;
            private int leeftijd;
            private String locatie;

            public Dier(String naam, int leeftijd, String locatie) {
                    this.naam = naam;
                    this.leeftijd = leeftijd;
                    this.locatie = locatie;
                }

            public void getInfo() {
                    System.out.println("Naam: " + naam + ", Leeftijd: " + leeftijd + ", Locatie: " + locatie);
                }
        }

    // Nieuwe klasse die Dier uitbreidt
    public class Hond extends Dier
        {

            private String ras;

            public Hond(String naam, int leeftijd, String locatie, String ras) {
                    super(naam, leeftijd, locatie); // roept de constructor van de superklasse Dier aan
                    this.ras = ras;
                }
            // je moet methodes niet opnieuw schrijven want een hond is ook een dier
            public void maakGeluid() {
                    System.out.println("Hond blaft");
                }

            // Herdefinieren (overriden): zelfde naam, zelfde parameters,
            // zelfde returntype als in Dier -> vervangt die van Dier.
            @Override
            public void getInfo() {
                    super.getInfo();          // eerst wat Dier al deed
                    System.out.println("Ras: " + ras);  // en dan het extra stuk
                }

            //alleen hond kan deze methode gebruiken
            public void speelFetch() {
                    System.out.println("Hond speelt fetch");
                }
        }

    // Gebruik van de klassen. Let op: code moet ALTIJD in een methode
    // staan, nooit los in de klasse.
    public class Dierentuin
        {
            private ArrayList<Dier> dieren;

            public Dierentuin() {
                    dieren = new ArrayList<Dier>();
                }

            public void vulAan() {
                    dieren.add(new Dier("Leeuw", 5, "Afrika"));
                    dieren.add(new Hond("Buddy", 3, "Thuis", "Labrador"));
                    // een Hond mag in een ArrayList<Dier>, want een Hond IS een Dier
                }

            public void toonAlles() {
                    for (Dier dier : dieren) {
                            dier.getInfo();   // Hond krijgt vanzelf zijn eigen versie
                            // dier.maakGeluid(); // FOUT: Dier kent die methode niet

                            if (dier instanceof Hond) {          // is het echt een Hond?
                                    Hond hond = (Hond) dier;         // casten
                                    hond.maakGeluid();
                                }
                        }
                }
        }
\end{codeblock}

\subsection{Polymorfie}

In \texttt{toonAlles()} staat \texttt{dier.getInfo()}. De variabele
\texttt{dier} is van het type \texttt{Dier}, maar er kan een \texttt{Hond} in
zitten. Java kijkt naar het echte object, niet naar het type van de variabele.
Zit er een \texttt{Hond} in, dan draait \texttt{getInfo()} van \texttt{Hond}. Je
hebt daar geen \texttt{if} voor nodig.

\begin{conceptbox}[Waarom is dat zo handig?]
    Zonder polymorfie moet je voor elk soort dier een \texttt{if} schrijven. Komt er
    later een \texttt{Kat} bij, dan moet je overal die \texttt{if}-ketens gaan
    aanpassen.

    \medskip
    M\'et polymorfie schrijf je \'e\'en lus. Een nieuwe subklasse zet er zijn eigen
    \texttt{@Override} bij en de lus werkt meteen mee. Dat heb je nodig bij een vraag
    als \textit{``bereken voor elke kandidaat het aantal juiste antwoorden''}, waarbij
    sommige kandidaten een joker mogen inzetten.
\end{conceptbox}

\begin{codeblock}[polymorfie in de praktijk]{java}
    // In de superklasse: de gewone regel
    public boolean heeftJuist(int vraag, char juist) {
            return antwoorden[vraag] == juist;
        }

    // In de subklasse: de uitzondering
    @Override
    public boolean heeftJuist(int vraag, char juist) {
            if (jokerVragen.contains(vraag)) {
                    return true;                        // joker = altijd juist
                }
            return super.heeftJuist(vraag, juist);  // rest ongewijzigd
        }

    // De teller hoeft van geen enkele joker te weten:
    public int aantalJuist(Kandidaat k) {
            int aantal = 0;
            for (int i = 0; i < 10; i++) {
                    if (k.heeftJuist(i, juisteAntwoorden[i])) { aantal++; }
                }
            return aantal;
        }
\end{codeblock}

\begin{examenbox}
    Zie je jezelf lange rijen \texttt{if (x instanceof Y)} schrijven, dan mis je
    waarschijnlijk een \texttt{@Override}. Gebruik \texttt{instanceof} enkel als de
    subklasse een methode heeft die de superklasse \'echt niet kent
    (zoals \texttt{speelFetch()}).
\end{examenbox}

\begin{center}
\begin{tikzpicture}[font=\small, x=1cm, y=-1cm]

\definecolor{superKleur}{RGB}{215,232,244}
\definecolor{subKleur}{RGB}{223,240,225}

\node[draw=black!55, fill=subKleur, minimum width=48mm, minimum height=10mm,
      align=left, font=\footnotesize\ttfamily, inner sep=3pt] at (0,0)
      {Hond-deel\\ ras};
\node[draw=black!55, fill=superKleur, minimum width=48mm, minimum height=13mm,
      align=left, font=\footnotesize\ttfamily, inner sep=3pt] at (0,1.2)
      {Dier-deel\\ naam, leeftijd, locatie};

\node[font=\footnotesize\bfseries, anchor=south] at (0,-0.7) {een Hond-object};

\draw[-{Stealth}, schoolBlue, thick] (-3.9, 1.2) -- (-2.55, 1.2);
\node[font=\tiny, schoolBlue, anchor=east, align=right] at (-4.0, 1.2)
      {1. \texttt{super(...)}\\ bouwt dit eerst};

\draw[-{Stealth}, schoolBlue, thick] (-3.9, 0) -- (-2.55, 0);
\node[font=\tiny, schoolBlue, anchor=east, align=right] at (-4.0, 0)
      {2. daarna vult de\\ Hond-constructor\\ de rest in};

\node[font=\tiny, anchor=west, align=left] at (2.7, 1.2)
      {\texttt{private} in Dier:\\ Hond kan er niet\\ rechtstreeks aan,\\ gebruik de getters};

\end{tikzpicture}
\end{center}

\begin{warningbox}[Regels bij overerving]
    \begin{itemize}
        \item \texttt{super(...)} moet de \textbf{eerste regel} van de constructor
              zijn.
        \item Een subklasse kan niet aan de \texttt{private} velden van de
              superklasse. Gebruik de getters.
        \item Overriden lukt alleen met \textbf{exact} dezelfde naam, parameters en
              returntype. Zet er \texttt{@Override} boven: typ je iets fout, dan
              zegt de compiler het meteen in plaats van dat je methode stil genegeerd
              wordt.
    \end{itemize}
\end{warningbox}
\section{Scanner}

\begin{examenbox}
    Op het examen krijg je de \textbf{Scanner-code} mee: de imports, het
    \texttt{try/catch}-blok en de regel \texttt{new Scanner(new File(...))}. Wat je
    \textbf{niet} krijgt, is welke methode je daarna moet oproepen. Dat je
    \texttt{next()} gebruikt voor een woord, \texttt{nextInt()} voor een geheel
    getal, \texttt{nextFloat()} of \texttt{nextDouble()} voor een kommagetal en
    \texttt{nextLine()} voor een hele regel --- en welke \texttt{hasNext\dots} daarbij
    hoort --- moet je zelf weten. Zet dat lijstje dus op je formularium.
\end{examenbox}

Een \texttt{Scanner} leest \textbf{token per token}. Spaties en regeleindes zijn
voor hem hetzelfde: allebei een scheiding. Het bestand is dus \'e\'en lange rij,
en een positie die opschuift bij elke \texttt{next\dots}

\begin{center}
\begin{tikzpicture}[
  tok/.style={draw=black!55, minimum height=6mm, font=\ttfamily\footnotesize,
              inner xsep=4pt, inner ysep=0pt},
  x=1cm, y=-1cm]

\definecolor{tokGelezen}{RGB}{235,238,242}

\node[tok, fill=tokGelezen] (t1) at (0,0)   {c};
\node[tok, fill=tokGelezen, anchor=west] (t2) at (0.45,0) {Minoes};
\node[tok, fill=white, anchor=west] (t3) at (2.1,0)  {4};
\node[tok, fill=white, anchor=west] (t4) at (2.6,0)  {d};
\node[tok, fill=white, anchor=west] (t5) at (3.1,0)  {Buddy};
\node[tok, fill=white, anchor=west] (t6) at (4.5,0)  {3};
\node[font=\footnotesize, anchor=west] at (5.1,0) {\dots};

\draw[-{Stealth}, schoolBlue, thick] (2.25,-0.55) -- (2.25,-0.2);
\node[font=\tiny, schoolBlue, anchor=north] at (2.25,-0.57) {positie};
\node[font=\tiny, anchor=east] at (-0.35,0) {gelezen};
\node[font=\tiny, anchor=west, align=left] at (5.6,0)
     {elke \texttt{next\dots}\\ schuift \'e\'en vakje op};

\end{tikzpicture}
\end{center}

Jij bepaalt met de volgorde van je oproepen hoe die rij ingedeeld wordt.

\begin{codeblock}[twee bronnen, dezelfde Scanner]{java}
    Scanner sc = new Scanner(System.in);            // toetsenbord
    Scanner sc = new Scanner(new File(naam));       // bestand
\end{codeblock}

\begin{codeblock}[scanner]{java}
    import java.util.Scanner;          // de Scanner zelf
    import java.io.File;               // om een bestand aan te duiden
    import java.io.FileNotFoundException; // de fout die je moet opvangen

    public class MijnKlasse
        {
            public MijnKlasse() {
                }

            public int sumOfTwoNumbers(int a, int b) {
                    return a + b;
                }

            // regel per regel lezen
            public void readFromFile(String filename) {
                    try {
                            Scanner sc = new Scanner(new File(filename));
                            while (sc.hasNextLine())
                            {
                                    String regel = sc.nextLine();
                                    System.out.println(regel);
                                }
                            sc.close();
                        }
                    catch (FileNotFoundException fnfe) {
                            System.out.println("file not found");
                        }
                }

            // getal per getal lezen, telkens twee tegelijk
            public void readNumbers(String filename) {
                    try {
                            Scanner sc = new Scanner(new File(filename));
                            while (sc.hasNextInt())
                            {
                                    int a = sc.nextInt();
                                    int b = sc.nextInt();
                                    System.out.println(sumOfTwoNumbers(a, b));
                                }
                            sc.close();
                        }
                    catch (FileNotFoundException fnfe) {
                            System.out.println("file not found");
                        }
                }
        }

    // belangrijke methodes van Scanner:
    // hasNext()      : is er nog een volgend token?
    // hasNextInt()   : en is dat token een geheel getal?
    // hasNextLine()  : is er nog een volgende regel?
    // next()         : lees het volgende token als String (1 woord)
    // nextInt()      : lees het volgende token als int
    // nextDouble()   : lees het volgende token als double
    // nextLine()     : lees de rest van de huidige regel als String
    // useDelimiter(";") : kies zelf het scheidingsteken
    // close()        : sluit de scanner (altijd doen)
\end{codeblock}

\subsubsection{Waarom die try en catch?}

\texttt{new Scanner(new File(naam))} kan mislukken: het bestand kan er gewoon
niet zijn. Java verplicht je daarom om dat geval af te handelen, anders
compileert je code niet (\textit{unreported exception FileNotFoundException}).

\begin{itemize}
    \item \texttt{try} --- ``probeer dit; het mag mislukken''.
    \item \texttt{catch (FileNotFoundException e)} --- ``en als het misloopt, doe
          dan dit in plaats van crashen''.
\end{itemize}

Heeft je methode een returntype, dan moet ze ook \textbf{in de catch} iets
teruggeven. Geef daar een duidelijke foutwaarde, meestal \texttt{-1} of
\texttt{null}.

\begin{warningbox}[De nextInt/nextLine-val]
    \texttt{nextInt()} leest het getal, maar laat de \textbf{enter} erachter staan.
    Roep je daarna \texttt{nextLine()} op, dan krijg je een lege String in plaats van
    de volgende regel. Oplossing: een extra \texttt{sc.nextLine()} inlassen om die
    enter op te souperen.
\end{warningbox}

\begin{conceptbox}[Welke hasNext gebruik je?]
    Laat je lusvoorwaarde passen bij wat je in de lus leest. Lees je
    \texttt{nextInt()}, gebruik dan \texttt{hasNextInt()}. Doe je dat niet en staat er
    tekst waar jij een getal verwacht, dan crasht je programma met
    \textit{InputMismatchException}.
\end{conceptbox}

\subsection{Een bestand met afwisselend tekst en getallen}

Per regel staan de gegevens van \'e\'en object, elk van een ander type. Stel dat
\texttt{dieren.txt} er zo uitziet:

\begin{codeblock}[dieren.txt]{}
    c Minoes 4
    d Buddy 3
    c Felix 7
    d Rex 11
\end{codeblock}

Per regel staan er drie tokens: een \texttt{char}, een \texttt{String} en een
\texttt{int}. Lees in de lus in \textbf{exact dezelfde volgorde} als in het
bestand.

\begin{codeblock}[afwisselend lezen]{java}
    public void leesDieren(String bestandsnaam)
    {
    try {
    Scanner sc = new Scanner(new File(bestandsnaam));

    while (sc.hasNext())    // is er nog een token?
    {
    String soort = sc.next();        // "c"  of "d"
    char type = soort.charAt(0);     // -> 'c' of 'd'
    String naam = sc.next();         // "Minoes"
    int leeftijd = sc.nextInt();     // 4

    Dier d = maakDier(type, naam, leeftijd);
    dieren.add(d);
    }
    sc.close();
    }
    catch (FileNotFoundException fnfe) {
            System.out.println("File not found");
        }
    }
\end{codeblock}

\begin{warningbox}[Er bestaat geen nextChar()]
    Wil je \'e\'en letter lezen, gebruik dan \texttt{sc.next()} en neem daarvan het
    eerste teken met \texttt{.charAt(0)}. Er is w\'el \texttt{nextInt()},
    \texttt{nextDouble()}, \texttt{nextFloat()}, \texttt{nextBoolean()} en
    \texttt{nextLong()}, maar geen \texttt{nextChar()}.
\end{warningbox}

\begin{warningbox}[Meng next() en nextLine() niet zomaar]
    Lees je een regel liever in \'e\'en keer in en splits je ze daarna zelf, doe dat
    dan \textbf{consequent}: \'ofwel alles met \texttt{nextLine()} en
    \texttt{split()}, \'ofwel alles token per token. Door elkaar gebruiken geeft de
    val die hierboven beschreven staat: \texttt{nextInt()} laat de enter staan en de
    \texttt{nextLine()} erna leest een lege String.
\end{warningbox}

Met een scheidingsteken in plaats van spaties, bijvoorbeeld
\texttt{Minoes;4;kat}, heb je twee opties:

\begin{codeblock}[bestand met puntkomma's]{java}
    // optie 1: regel voor regel en zelf splitsen
    while (sc.hasNextLine()) {
            String regel = sc.nextLine();          // "Minoes;4;kat"
            String[] delen = regel.split(";");
            String naam = delen[0];                // "Minoes"
            int leeftijd = Integer.parseInt(delen[1]);  // 4
            String soort = delen[2];               // "kat"
        }

    // optie 2: de Scanner zelf op de puntkomma laten splitsen
    sc.useDelimiter("\\s*;\\s*");   // ; met eventueel spaties errond
    String naam = sc.next();
    int leeftijd = sc.nextInt();
\end{codeblock}

\begin{examenbox}
    Optie 1 is in de praktijk veiliger: je ziet meteen wat er in \texttt{delen[]}
    zit en je kan met \texttt{delen.length} controleren of de regel wel volledig is.
    Vergeet niet dat \texttt{split()} altijd \texttt{String}'s teruggeeft --- getallen
    moet je zelf omzetten met \texttt{Integer.parseInt()}.
\end{examenbox}

\section{Random}

\texttt{Random} maakt willekeurige getallen. Importeer
\texttt{java.util.Random}.

\begin{codeblock}[random]{java}
    import java.util.Random; //importeer de Random klasse

    public class MijnKlasse
        {

            public MijnKlasse() {
                }

            public void generateRandomNumbers() {
                    Random rand = new Random(); // maak een Random object

                    // genereer een willekeurig geheel getal tussen 0 (inclusief) en 100 (exclusief)
                    int randomInt = rand.nextInt(100);
                    System.out.println("Willekeurig geheel getal: " + randomInt);

                    // genereer een willekeurig decimaal getal tussen 0.0 (inclusief) en 1.0 (exclusief)
                    float randomFloat = rand.nextFloat();
                    System.out.println("Willekeurig decimaal getal: " + randomFloat);

                    // een willekeurig getal binnen een bereik, bv. 50 t.e.m. 150
                    int min = 50;
                    int max = 150;
                    int randomInRange = rand.nextInt(max - min + 1) + min;
                    // de +1 is nodig omdat nextInt(n) stopt bij n-1
                    System.out.println("Getal tussen " + min + " en " + max + ": " + randomInRange);
                }
        }
\end{codeblock}

\begin{warningbox}[nextInt(n) is exclusief]
    \texttt{nextInt(10)} geeft 0 t.e.m.\ \textbf{9}, nooit 10. Voor een dobbelsteen
    schrijf je dus \texttt{rand.nextInt(6) + 1}. Voor het bereik
    \texttt{min} t.e.m.\ \texttt{max} is de formule
    \texttt{rand.nextInt(max - min + 1) + min}.
\end{warningbox}

\begin{examenbox}
    Maak je \texttt{Random}-object \'e\'en keer aan, als veld van je klasse, niet
    telkens opnieuw binnen een lus.
\end{examenbox}

\section{UML diagrammen}
\label{sec:uml}

Het eerste deel van het examen is een \keyterm{UML-klassendiagram} op papier,
goed voor \textbf{20\% van je punten}. Je krijgt een opgave in gewone taal. Daaruit
haal je de klassen, hun attributen, de signatuur van hun methodes en de relaties
ertussen. Je schrijft geen code.

\subsection{Hoe ziet \'e\'en klasse eruit?}

Een klasse is een rechthoek met \textbf{drie vakken}: de naam, de attributen en
de methodes.

\begin{center}
    \begin{tikzpicture}[font=\small\ttfamily]
        \node[rectangle split, rectangle split parts=3, draw, thick,
            rectangle split part align=left, text width=7.4cm,
            rectangle split part fill={schoolBlue!15, white, white}] (k)
        {
            \textbf{\normalfont Kandidaat}
            \nodepart{two}
            - mobielNummer : String \\
            - geboortejaar : int \\
            - antwoorden : char[] \\
            \underline{+ AANTAL\_VRAGEN : int}
            \nodepart{three}
            + Kandidaat(nr : String, jaar : int) \\
            + getMobielNummer() : String \\
            + heeftJuist(i : int, c : char) : boolean \\
            - isGeldigNummer() : boolean
        };
    \end{tikzpicture}
\end{center}

\begin{center}
    \begin{tabular}{@{}clp{7.4cm}@{}}
        \toprule
        \textbf{Teken}         & \textbf{Betekenis} & \textbf{Wanneer}                              \\
        \midrule
        \texttt{-}             & private            & standaard voor \textbf{alle} attributen       \\
        \texttt{+}             & public             & standaard voor methodes die anderen gebruiken \\
        \texttt{\#}            & protected          & zelden nodig in dit vak                       \\
        \midrule
        \underline{onderlijnd} & static             & constanten, tellers voor alle objecten        \\
        \textit{schuin}        & abstract           & klasse of methode zonder invulling            \\
        \bottomrule
    \end{tabular}
\end{center}

\begin{warningbox}[Hier gaan de punten verloren]
    Schrijf \textbf{altijd} het type mee, in de vorm \texttt{naam : type}. En bij
    methodes de \textbf{volledige signatuur}: elke parameter met zijn type, plus het
    returntype na de dubbele punt. Een methode zonder returntype schrijf je als
    \texttt{: void}. Enkel ``\texttt{berekenScore()}'' neerschrijven levert je maar
    een deel van de punten op.
\end{warningbox}

\subsection{De relaties tussen klassen}

Je hebt drie relaties nodig. Vergeet ze niet te tekenen: ze staan op punten.

\begin{center}
    \begin{tikzpicture}[font=\small\ttfamily, node distance=2.2cm,
            every node/.style={align=center}]
        % overerving
        \node[draw, thick, minimum width=2.4cm, minimum height=0.8cm] (sup) {Kandidaat};
        \node[draw, thick, minimum width=2.4cm, minimum height=0.8cm, below=1.6cm of sup] (sub) {VIPKandidaat};
        \draw[-{Triangle[open, length=3.5mm, width=3.5mm]}, thick] (sub) -- (sup);
        \node[right=0.15cm of sub, font=\small\sffamily, anchor=west, xshift=0.6cm]
        {\textit{is een} (extends)};

        % aggregatie
        \node[draw, thick, minimum width=2.4cm, minimum height=0.8cm, right=5.2cm of sup] (quiz) {Quiz};
        \node[draw, thick, minimum width=2.4cm, minimum height=0.8cm, below=1.6cm of quiz] (kand) {Kandidaat};
        \draw[{Diamond[open, length=4mm, width=2.6mm]}-, thick] (quiz) -- (kand)
        node[midway, right, font=\scriptsize\sffamily] {\quad heeft};
        \node[left=0.05cm of quiz.south west, font=\scriptsize, anchor=north east, yshift=-0.1cm] {1};
        \node[left=0.05cm of kand.north west, font=\scriptsize, anchor=south east, yshift=0.1cm] {0..*};
    \end{tikzpicture}
\end{center}

\begin{enumerate}
    \item \keyterm{Overerving} --- pijl met een \textbf{open driehoek} van de
          subklasse naar de superklasse. Lees ze als ``is een''. De geerfde
          velden herhaal je \textbf{niet} in de subklasse: die staan al boven.
    \item \keyterm{Aggregatie / compositie} --- lijn met een \textbf{ruit} aan de
          kant van de klasse die de collectie bezit. Lees ze als ``heeft''.
          Hier hoort de \keyterm{multipliciteit} bij: \texttt{1} bij de Quiz en
          \texttt{0..*} bij de Kandidaat, want \'e\'en quiz heeft nul of meer
          kandidaten.
    \item \keyterm{Gebruikt} --- gewone pijl met open punt. Klasse A gebruikt B,
          maar bezit ze niet: B komt enkel voor als parametertype of als lokale
          variabele.
\end{enumerate}

\begin{center}
    \begin{tabular}{@{}lp{7.4cm}@{}}
        \toprule
        \textbf{Multipliciteit} & \textbf{Betekenis}                                         \\
        \midrule
        \texttt{1}              & precies \'e\'en                                            \\
        \texttt{0..1}           & nul of \'e\'en (mag \texttt{null} zijn)                    \\
        \texttt{0..*}           & nul of meer --- de gewone keuze bij een \texttt{ArrayList} \\
        \texttt{1..*}           & minstens \'e\'en                                           \\
        \texttt{10}             & precies tien (bv.\ een \texttt{char[10]} met antwoorden)   \\
        \bottomrule
    \end{tabular}
\end{center}

\subsection{Van tekst naar diagram: de werkwijze}

Ga zin per zin door de opgave.

\begin{enumerate}
    \item \textbf{Zelfstandige naamwoorden} $\rightarrow$ kandidaat-klassen.
          ``Quiz'', ``kandidaat'', ``ticket'', ``mier''. Niet elk naamwoord wordt
          een klasse: iets wordt pas een klasse als het \emph{eigen gegevens
              \'en eigen gedrag} heeft.
    \item \textbf{``wordt gekenmerkt door'', ``moet ingeven''} $\rightarrow$
          attributen. Kies meteen het type, en denk na:
          een telefoonnummer is een \texttt{String} (anders valt de \texttt{0}
          van \texttt{04\dots} weg), een antwoord \texttt{'A'} t.e.m. \texttt{'D'}
          is een \texttt{char}, een schiftingsantwoord is een \texttt{double}.
    \item \textbf{``heeft meerdere'', ``een lijst van''} $\rightarrow$ een
          collectie, dus een ruit-relatie met \texttt{0..*}.
    \item \textbf{``is ook een \dots\ maar met extra''} $\rightarrow$ overerving.
          Test het met de zin ``een X \emph{is een} Y''. Klopt dat niet, gebruik
          dan geen \texttt{extends}.
    \item \textbf{Werkwoorden} $\rightarrow$ methodes. Het returntype staat
          bijna letterlijk in de opgave:
\end{enumerate}

\begin{center}
    \begin{tabular}{@{}p{6.2cm}l@{}}
        \toprule
        \textbf{Zin in de opgave}                                        & \textbf{Returntype}          \\
        \midrule
        ``wordt gecontroleerd of \dots, anders wordt hij niet aanvaard'' & \texttt{boolean}             \\
        ``het aantal verwijderde \dots\ wordt teruggegeven''             & \texttt{int}                 \\
        ``geeft een collectie met deze \dots\ terug''                    & \texttt{ArrayList<X>}        \\
        ``de winnaar aanduiden''                                         & \texttt{X} (het object zelf) \\
        ``bereken het gemiddelde''                                       & \texttt{double}              \\
        ``toon / print''                                                 & \texttt{void}                \\
        \bottomrule
    \end{tabular}
\end{center}

\textbf{Checklist voor je je blad afgeeft}
\begin{enumerate}
    \item Staat bij \'elk attribuut en \'elke methode een \texttt{+} of \texttt{-}?
    \item Staat overal het type?
    \item Heb je de lijn tussen de klassen \'echt getekend, met ruit \'en
          multipliciteit?
    \item Staat de constructor bij de methodes?
    \item Heb je geerfde velden niet dubbel geschreven?
\end{enumerate}


\section{HashMaps}
\label{sec:hashmap}

\subsection{Wat is een HashMap?}

Denk aan een \textbf{array waarbij jij de indexen zelf kiest}.

Bij een gewone array zijn de plaatsen genummerd: \texttt{namen[0]},
\texttt{namen[1]}, \texttt{namen[2]}. Je m\'o\'et een geheel getal gebruiken, en
je moet dus zelf onthouden dat Bob op plaats 3 staat.

Bij een \texttt{HashMap} mag die ``index'' van eender welk type zijn: een
\texttt{String}, een \texttt{char}, of zelfs een object van je eigen klasse. Je
schrijft dan gewoon \texttt{mieren.get("A7")} in plaats van
\texttt{mieren[3]}.

\begin{center}
    \begin{tabular}{@{}lll@{}}
        \toprule
                             & \textbf{array}          & \textbf{HashMap}             \\
        \midrule
        Plaats aanduiden met & \texttt{0, 1, 2, \dots} & wat jij wil: \texttt{"A7"}   \\
        Schrijven            & \texttt{a[3] = ant;}    & \texttt{map.put("A7", ant);} \\
        Lezen                & \texttt{a[3]}           & \texttt{map.get("A7")}       \\
        Lengte               & vast bij het aanmaken   & groeit vanzelf               \\
        Volgorde             & vast                    & \textbf{geen} volgorde       \\
        \bottomrule
    \end{tabular}
\end{center}

Die zelfgekozen index heet de \keyterm{sleutel} (\textit{key}), en wat erachter
zit heet de \keyterm{waarde} (\textit{value}). Een \texttt{HashMap} is dus een
verzameling \textbf{paren}: bij elke sleutel hoort \'e\'en waarde. Denk aan een
woordenboek: je zoekt op het woord (sleutel) en krijgt de vertaling (waarde).

\subsection{Twee dingen waar je ze voor gebruikt}

\textbf{1. Iets terugvinden zonder te zoeken.} Je hebt een id en je wil het
object dat erbij hoort.

\medskip
\textbf{2. Bijhouden hoe vaak iets voorkomt.} Je weet niet op voorhand welke
waarden je gaat tegenkomen. Met een array lukt dat niet: daarvoor moet je vooraf
al een vakje per mogelijk woord klaarzetten. Een \texttt{HashMap} maakt het vakje
aan zodra je het woord voor het eerst ziet, en telt er daarna bij.

\begin{codeblock}[komt een waarde meerdere keren voor?]{java}
    // sleutel = het woord, waarde = hoe vaak we het al zagen
    HashMap<String, Integer> tellers = new HashMap<>();

    tellers.put("bob", 1);   // eerste keer bob
    tellers.put("an", 1);    // eerste keer an
    tellers.put("bob", 2);   // bob voor de tweede keer

    tellers.get("bob")            // 2   -> komt meermaals voor
    tellers.get("an")             // 1   -> komt maar een keer voor
    tellers.get("zoe")            // null-> nog nooit gezien
    tellers.containsKey("bob")    // true
\end{codeblock}

De volledige uitwerking staat verderop bij \textit{Tellen hoe vaak iets
    voorkomt}.

\subsection{Waarom bestaat dit naast ArrayList?}

Stel dat je een mier moet terugvinden op basis van haar id. Met een
\texttt{ArrayList} doe je dat zo:

\begin{codeblock}[zoeken in een ArrayList]{java}
    public Ant zoek(String id) {
            for (Ant a : ants) {                    // je overloopt ALLES
                    if (a.getId().equals(id)) { return a; }
                }
            return null;
        }
\end{codeblock}

Bij 1000 mieren kijkt die lus in het slechtste geval 1000 keer. Een
\texttt{HashMap} rekent de plaats van de sleutel rechtstreeks uit en vindt het
antwoord in \'e\'en stap, hoe groot de verzameling ook is:

\begin{codeblock}[zoeken in een HashMap]{java}
    Ant a = ants.get(id);   // klaar, geen lus
\end{codeblock}

\begin{conceptbox}[Wanneer kies je wat?]
    \textbf{ArrayList} als de \textbf{volgorde} telt of als je gewoon ``alles
    overlopen'' nodig hebt: een lijst deelnemers, een lijst tickets.

    \medskip
    \textbf{HashMap} als je iets steeds \textbf{opzoekt via een unieke sleutel}:
    mier per pheromone-id, student per studentennummer, reservatie per attractie.

    \medskip
    \textbf{HashSet} als je enkel wil weten \textbf{of iets er al is}, zonder
    bijhorende waarde: reeds gebruikte codes, unieke id's. Een \texttt{HashSet} is in
    feite een \texttt{HashMap} zonder waarden; dubbels worden stil genegeerd.
\end{conceptbox}

\subsection{Syntax}

Je schrijft twee types tussen de punthaken: eerst die van de sleutel, dan die van
de waarde.

\begin{codeblock}[hashmap basis]{java}
    import java.util.HashMap;

    public class Mierenhoop
        {
            // sleutel = String (het id), waarde = Ant (het object)
            private HashMap<String, Ant> mieren;

            public Mierenhoop() {
                    mieren = new HashMap<String, Ant>(); // of new HashMap<>();
                }

            public void voegToe(Ant a) {
                    mieren.put(a.getId(), a);   // (sleutel, waarde)
                }

            public Ant zoek(String id) {
                    return mieren.get(id);      // null als de sleutel niet bestaat
                }
        }
\end{codeblock}

\begin{center}
    \begin{tabular}{@{}llp{5.9cm}@{}}
        \toprule
        \textbf{Methode}          & \textbf{Geeft terug}    & \textbf{Doet}                   \\
        \midrule
        \texttt{put(k, v)}        & vorige waarde           & paar toevoegen of overschrijven \\
        \texttt{get(k)}           & waarde of \texttt{null} & waarde bij die sleutel          \\
        \texttt{containsKey(k)}   & \texttt{boolean}        & bestaat die sleutel?            \\
        \texttt{containsValue(v)} & \texttt{boolean}        & komt die waarde voor?           \\
        \texttt{remove(k)}        & verwijderde waarde      & paar verwijderen                \\
        \texttt{size()}           & \texttt{int}            & aantal paren                    \\
        \texttt{isEmpty()}        & \texttt{boolean}        & leeg?                           \\
        \texttt{keySet()}         & \texttt{Set<K>}         & alle sleutels                   \\
        \texttt{values()}         & \texttt{Collection<V>}  & alle waarden                    \\
        \bottomrule
    \end{tabular}
\end{center}

\begin{warningbox}[put overschrijft stil]
    Bestaat de sleutel al, dan \textbf{vervangt} \texttt{put} de oude waarde zonder
    waarschuwing. Wil je dat niet, check dan eerst met \texttt{containsKey}. Een
    sleutel kan dus maar \'e\'en keer voorkomen; een waarde mag wel meermaals
    voorkomen.
\end{warningbox}

\subsection{Een HashMap doorlopen}

Er zijn drie manieren, afhankelijk van wat je nodig hebt.

\begin{codeblock}[hashmap overlopen]{java}
    // 1. enkel de sleutels
    for (String id : mieren.keySet()) {
            System.out.println(id);
        }

    // 2. enkel de waarden (het vaakst gebruikt)
    for (Ant a : mieren.values()) {
            totaal += a.getLoad();
        }

    // 3. sleutel en waarde samen
    for (String id : mieren.keySet()) {
            Ant a = mieren.get(id);
            System.out.println(id + " draagt " + a.getLoad());
        }
\end{codeblock}

\begin{warningbox}[Een HashMap heeft geen volgorde]
    De volgorde waarin \texttt{keySet()} of \texttt{values()} de elementen teruggeeft
    ligt \textbf{niet} vast, en is niet de volgorde waarin je ze toevoegde. Heb je
    volgorde nodig, gebruik dan een \texttt{ArrayList}. Verwijderen tijdens het
    overlopen mag hier evenmin: ook een \texttt{HashMap} geeft dan een
    \textit{ConcurrentModificationException}.
\end{warningbox}

\subsection{Tellen hoe vaak iets voorkomt}

Hoe vaak komt elk woord voor? Sleutel is het woord, waarde is de teller.

\begin{codeblock}[tellen met een HashMap]{java}
    public HashMap<String, Integer> telWoorden(ArrayList<String> woorden)
    {
            HashMap<String, Integer> tellers = new HashMap<>();
            for (String w : woorden)
            {
                    if (tellers.containsKey(w)) {
                            tellers.put(w, tellers.get(w) + 1); // ophogen
                        } else {
                            tellers.put(w, 1);                  // eerste keer
                        }
                }
            return tellers;
        }
\end{codeblock}

\begin{warningbox}[Geen int, wel Integer]
    Net als bij \texttt{ArrayList} mag je in de punthaken geen basistype schrijven.
    Het is \texttt{HashMap<String, Integer>}, niet
    \texttt{HashMap<String, int>}. Voor kommagetallen gebruik je \texttt{Double},
    voor \texttt{true}/\texttt{false} \texttt{Boolean}.
\end{warningbox}

\subsection{HashSet}

\begin{codeblock}[hashset]{java}
    import java.util.HashSet;

    private HashSet<String> gebruikteCodes = new HashSet<>();

    public boolean voerCodeIn(String code) {
            if (gebruikteCodes.contains(code)) {
                    return false;         // code al eens gebruikt
                }
            gebruikteCodes.add(code); // add geeft false als ze er al in zat
            return true;
        }
\end{codeblock}

\begin{examenbox}
    Staat het woord \textbf{``uniek''} in de opgave --- een uniek id, een code die
    je maar \'e\'en keer mag gebruiken --- gebruik dan een \texttt{HashSet} of de
    sleutels van een \texttt{HashMap}. Je hebt dan geen zoeklus nodig.
\end{examenbox}


\section{Compilerfouten lezen}

De compiler zegt wat er fout is, maar in zijn eigen woorden. Deze tabel
vertaalt de boodschappen die je in dit vak tegenkomt.

\begin{center}
    \begin{tabular}{@{}p{5.2cm}p{9.2cm}@{}}
        \toprule
        \textbf{Boodschap}                                                       & \textbf{Wat er echt aan de hand is} \\
        \midrule
        \texttt{cannot find symbol: variable x}                                  &
        Er bestaat geen veld of variabele met die naam. Typfout, of je schrijft
        \texttt{this.x} terwijl je veld anders heet.                                                                   \\
        \addlinespace
        \texttt{variable x might not have been initialized}                      &
        Een lokale variabele zonder startwaarde. Velden krijgen automatisch 0 /
        \texttt{false} / \texttt{null}, variabelen in een methode niet.                                                \\
        \addlinespace
        \texttt{not a statement}                                                 &
        Een regel die een waarde oplevert maar nergens naartoe gaat, bv.
        \texttt{rn.nextInt(6) + 1;}. Zet er \texttt{int d =} voor.                                                     \\
        \addlinespace
        \texttt{missing return statement}                                        &
        Er is een pad door je methode dat niets teruggeeft. Meestal een \texttt{if}
        zonder \texttt{else}.                                                                                          \\
        \addlinespace
        \texttt{incompatible types: possible lossy conversion from float to int} &
        Je steekt een kommagetal in een \texttt{int}. Kies het juiste type, ook voor het
        returntype van de methode.                                                                                     \\
        \addlinespace
        \texttt{constructor X in class X cannot be applied to given types}       &
        De superklasse heeft geen constructor zonder parameters, dus moet jij
        \texttt{super(...)} zelf schrijven.                                                                            \\
        \addlinespace
        \texttt{invalid method declaration; return type required}                &
        Je dacht een constructor te schrijven, maar de naam wijkt af van de klasse ---
        let op hoofdletters: \texttt{Antfarm} $\neq$ \texttt{AntFarm}.                                                 \\
        \addlinespace
        \texttt{class, interface, or enum expected}                              &
        Accolades kloppen niet, of er staat een losse \texttt{*/} in je bestand.                                       \\
        \bottomrule
    \end{tabular}
\end{center}

\subsection{Fouten die pas bij het uitvoeren komen}

\begin{center}
    \begin{tabular}{@{}p{5.2cm}p{9.2cm}@{}}
        \toprule
        \textbf{Boodschap}                            & \textbf{Wat er echt aan de hand is} \\
        \midrule
        \texttt{NullPointerException: \dots\ is null} &
        Je array of \texttt{ArrayList} is nooit aangemaakt. Zet
        \texttt{lijst = new ArrayList<>();} in de constructor.                              \\
        \addlinespace
        \texttt{ArrayIndexOutOfBounds\-Exception}     &
        Je schrijft buiten de array. Meestal \texttt{<=} waar \texttt{<} moest staan, of
        je telt de index op v\'o\'or je opslaat.                                            \\
        \addlinespace
        \texttt{ConcurrentModification\-Exception}    &
        Je verwijdert uit een lijst terwijl je er met een for-each over loopt. Gebruik
        een \texttt{Iterator} of een hulplijst.                                             \\
        \bottomrule
    \end{tabular}
\end{center}

\subsection{Veld of lokale variabele?}

De compiler geeft hier geen fout, want de code is geldig. Ze crasht pas bij
het uitvoeren.

\begin{codeblock}[FOUT: het veld blijft null]{java}
    public class AntFarm
        {
            private ArrayList<Ant> ants;

            public AntFarm()
            {
                    ArrayList<Ant> ants = new ArrayList<Ant>();
                    // dit is een NIEUWE doos die alleen hier bestaat.
                    // Zodra de constructor klaar is, is ze weg.
                    // Het veld hierboven staat nog altijd op null.
                }
        }
\end{codeblock}

\begin{codeblock}[JUIST]{java}
    public AntFarm()
    {
            ants = new ArrayList<Ant>();       // vult het veld
            // of, nog duidelijker:
            this.ants = new ArrayList<Ant>();
        }
\end{codeblock}

\begin{conceptbox}[De regel in \'e\'en zin]
    \textbf{Type ervoor $=$ nieuwe variabele maken. Geen type ervoor $=$ bestaande
        variabele invullen.} In een constructor wil je altijd het tweede.
\end{conceptbox}

Schrijf altijd \texttt{this.} in je constructor. Typ je dan per ongeluk toch een
type, dan krijg je meteen een compilerfout in plaats van een verrassing later ---
\texttt{this.ArrayList<Ant> ants} bestaat namelijk niet.

\begin{warningbox}[Zelfde probleem, andere vorm]
    \begin{itemize}
        \item \texttt{return ArrayList<Ant> ants;} --- bij \texttt{return} zet je nooit
              een type. Enkel \texttt{return ants;}.
        \item \texttt{return [];} --- Java kent geen \texttt{[]} voor een lege lijst,
              dat is Python. Gebruik \texttt{new ArrayList<Ant>()}.
        \item \texttt{for (i = 0; i < n; i++)} --- het type hoort binnen de haakjes:
              \texttt{for (int i = 0; \dots)}.
    \end{itemize}
\end{warningbox}

\section{Trucs die telkens terugkomen}

\subsection{De while-lus die zichzelf corrigeert}

Moet je iets willekeurigs trekken dat aan een voorwaarde voldoet, tel dan niet
zelf. Laat de lus gewoon opnieuw proberen.

\begin{codeblock}[8 tekens, elk teken maximaal twee keer]{java}
    public String generatePheromoneId()
    {
            String charToUse = "0123456789aeiou";
            String returnId = "";
            Random rn = new Random();

            while (returnId.length() < 8)
            {
                    int index = rn.nextInt(charToUse.length());
                    char c = charToUse.charAt(index);

                    if (aantalKeer(returnId, c) < 2)
                    {
                            returnId += c;
                        }
                    // past het niet, dan groeit returnId niet
                    // en probeert de while vanzelf opnieuw
                }
            return returnId;
        }
\end{codeblock}

\keyterm{Waarom dit werkt:} de voorwaarde van de \texttt{while} kijkt naar het
\emph{resultaat}, niet naar het aantal pogingen. Voeg je niets toe, dan is de
voorwaarde nog altijd waar en draait de lus nog eens. Geen teller die je moet
terugdraaien met \texttt{i--}.

\subsection{Een hulpmethode voor het telwerk}

Zodra je een lus in een lus krijgt, zet de binnenste in een eigen
\texttt{private} methode. Je hoofdmethode blijft leesbaar en je kan het stuk
apart testen.

\begin{codeblock}[private hulpmethode]{java}
    private int aantalKeer(String tekst, char c)
    {
    int aantal = 0;
    for (int i = 0; i < tekst.length(); i++)
    {
    if (tekst.charAt(i) == c)   // bij char mag ==
    {
    aantal++;
    }
    }
    return aantal;
    }
\end{codeblock}

\texttt{private} omdat alleen deze klasse hem gebruikt. Zet je hem
\texttt{public}, dan verwacht de lezer dat het een echte functie van de klasse is.

\subsection{Een teken tellen zonder hulpmethode}

\begin{codeblock}[indexOf en lastIndexOf]{java}
    if (returnId.indexOf(c) == returnId.lastIndexOf(c))
    {
            // c zit er 0 of 1 keer in
        }
\end{codeblock}

Zit het teken er niet in, dan geven beide \texttt{-1}. Zit het er \'e\'en keer
in, dan geven beide dezelfde plaats. Pas vanaf twee keer verschillen ze.

\subsection{Een willekeurig teken uit een tekst}

\begin{codeblock}[char uit een String]{java}
    String s = "abcd";
    char c = s.charAt(rn.nextInt(s.length()));
\end{codeblock}

Dit is hetzelfde patroon als een willekeurig element uit een lijst:
\texttt{lijst.get(rn.nextInt(lijst.size()))}.

\subsection{Startwaarden van een teller}

\begin{center}
    \begin{tabular}{@{}lll@{}}
        \toprule
        \textbf{Wat je doet} & \textbf{Start op}             & \textbf{Waarom}              \\
        \midrule
        optellen             & \texttt{0}                    & 0 verandert een som niet     \\
        vermenigvuldigen     & \texttt{1}                    & 1 verandert een product niet \\
        grootste zoeken      & eerste element of \texttt{-1} & alles is groter dan dat      \\
        kleinste zoeken      & eerste element                & anders vind je nooit iets    \\
        object zoeken        & \texttt{null}                 & ``niets gevonden''           \\
        \bottomrule
    \end{tabular}
\end{center}

Zoek je de grootste, start dan liever op \texttt{lijst.get(0)} dan op een zelf
gekozen getal. Met negatieve waarden in de lijst gaat \texttt{0} namelijk fout.

\subsection{Eerst doen, dan pas controleren}

Een testklasse verklapt vaak de volgorde. Kijk naar deze drie regels:

\begin{codeblock}[uit AntFarmTest]{java}
    assertFalse( a1.carryLoad( 5.0f ) );                     // geeft false
    assertEquals( 13.8f, a1.totalWeightCarried(), 0.001f );  // maar telt wel mee
    assertEquals( 5, a1.amountLoadsCarried() );
\end{codeblock}

\texttt{false} betekent hier niet ``er is niets gebeurd'', maar ``dit was de
laatste''. De lading wordt dus \textbf{eerst opgeslagen} en pas daarna kijk je of
de mier sterft.

\begin{codeblock}[de juiste volgorde]{java}
    public boolean carryLoad(float weight)
    {
            if (dead)
            {
                    return false;              // dode mier: niets doen
                }

            loadsCarried[loadIndex] = weight;   // 1. opslaan
            loadIndex++;                        // 2. index verhogen

            if (totalWeightCarried() >= 50 || loadIndex >= 5)   // 3. pas nu testen
                {
                    die();
                    return false;
                }
            return true;
        }
\end{codeblock}

\begin{examenbox}
    Lees de testcode als de opgave. Elke \texttt{assert} zegt precies wat er moet
    gebeuren, ook in welke volgorde. Twijfel je tussen \texttt{>} en \texttt{>=},
    zoek dan de test met de exacte grenswaarde --- hier 49.9 $+$ 0.1 $=$ 50.0, dus
    \texttt{>=}.
\end{examenbox}

\subsection{Opslaan en dan pas ophogen}

\begin{codeblock}[FOUT en JUIST]{java}
    // FOUT: plaats 0 blijft leeg, je schrijft vanaf plaats 1
    loadIndex++;
    loadsCarried[loadIndex] = weight;

    // JUIST
    loadsCarried[loadIndex] = weight;
    loadIndex++;
\end{codeblock}

\texttt{loadIndex} wijst altijd naar de \textbf{eerstvolgende vrije plaats}.
Daarom is hij tegelijk het aantal opgeslagen elementen, en kan je hem gebruiken
als grens in je lus: \texttt{for (int i = 0; i < loadIndex; i++)}.

\subsection{Een array precies groot genoeg maken}

\begin{codeblock}[kopie zonder lege plaatsen]{java}
    public float[] getExactLoadsCarried()
    {
            float[] kopie = new float[loadIndex];   // exact zo groot als nodig
            for (int i = 0; i < loadIndex; i++)
            {
                    kopie[i] = loadsCarried[i];
                }
            return kopie;
        }
\end{codeblock}

Ook hier geen for-each: je hebt de index nodig om in de nieuwe array te
schrijven.

\end{document}
